\section{Introduction}
Let $G = (V,w)$ be a weighted graph, where $w : \binom{V}{2} \rightarrow \R_{\ge 0}$ assigns a non-negative weight to every edge slot.  
We denote the edges of $G$, i.e.\ the edge slots that are given positive weight, by $E(G)$.  
For a nontrivial set $\emptyset \ne X \subsetneq V$ let 
$\Delta_G(X)$ be the set of edges of $G$ with exactly one endpoint in $X$ and 
one endpoint in $\overline{X} = V \setminus X$.  A \emph{cut} of $G$ is a set of 
edges of the form $\Delta_G(X)$ for some nontrivial set $X \subseteq V$.  We call $X$ and 
$\overline{X}$ the \emph{shores} of the cut.  The minimum cut problem is to determine the minimum of 
$\sum_{e \in \Delta_G(X)} w(e)$ over all non trivial subsets $X$.    This is equivalent to the minimum total weight of edges that need to be removed 
from $G$ in order to disconnect it.  We call this minimum value  $\lambda(G)$.  A set of edges $\Delta_G(X)$ realizing $\lambda(G)$ 
is called a \emph{minimum cut} of $G$.  If $G$ is unweighted $\lambda(G)$ is known as the \emph{edge connectivity} of $G$ and is the 
minimum number of edges whose removal disconnects $G$.

Computing the weight of a minimum cut of a graph is a 
fundamental computational problem that has been extensively studied in theoretical computer 
science since at least the 1960s \cite{GH61, FF62}.  It is also a problem of great practical importance, with applications 
to clustering algorithms \cite{Bot93} and evaluating network reliability, among others (see \cite{PQ82} for a survey of applications).  
Classically it is known that edge connectivity can be computed in nearly linear time even by 
\emph{deterministic} algorithms \cite{KT19,HRW20}.  For weighted graphs with $m$ edges, the weight of a minimum cut can be determined in 
nearly linear time\footnote{The $\tO(\cdot)$ notation hides polylogarithmic factors in its argument.} $\tO(m)$ by a randomized algorithm \cite{Karger00,MN20,GMW20} and in almost linear time $O(m^{1+o(1)})$ by a deterministic algorithm \cite{Li20}.

In this work we study quantum algorithms for the minimum cut problem
in two standard models for graph problems, the adjacency matrix and 
the adjacency array models.  In the adjacency matrix model a query consists of a pair $\{u,v\}$ of vertices, and the answer is $w(\{u,v\})$.  The adjacency array model allows $3$ types of queries: 
one can query the degree of a vertex $v$, the name of the $\ith$ neighbor of $v$, according to some arbitrary ordering, and the weight of the edge between 
$v$ and its $\ith$ neighbor.

For classical randomized algorithms, in the adjacency matrix model it is known that even deciding if a graph is connected or not requires $\Omega(n^2)$ queries in the worst case \cite{DHHM06}.
More recently, the randomized query complexity of edge connectivity was studied by Bishnu, Ghosh, Mishra and Paraashar \cite{BGMP20} in a common generalization of the adjacency matrix 
and adjacency array models called the \emph{local query} model.  This model allows queries to the degree of a vertex and to the $\ith$ neighbor of a vertex $v$, 
as in the adjacency array model, and also queries as to whether or not $\{u,v\}$ is an edge, as in the adjacency matrix model.  Over simple graphs $G$ with $m$ edges, they show an 
$\Omega(m)$ lower bound on the number of local queries needed by a randomized algorithm to succeed with probability $2/3$ for both the problems of determining the edge connectivity and 
outputting a cut realizing the edge connectivity \cite[Theorems 2 and 3]{BGMP20}.  

In this work we completely characterize the quantum query and time complexity of the minimum cut problem in the adjacency matrix model.  The complexity depends on what we call 
the \emph{edge-weight ratio}.  We say a graph has edge-weight ratio $\tau$ if the ratio of the largest weight of the graph to the smallest is at most $\tau$.
When the edge-weight ratio of an $n$-vertex graph is $\tau$, we give a bounded-error quantum algorithm to solve the minimum cut problem using $\tilde O(n^{3/2} \sqrt{\tau})$ queries and time 
in the adjacency matrix model (\cref{thm:time}).
For the unweighted case, i.e.\ the case $\tau = 1$, one can see this bound is tight as D\"{u}rr, Heiligman, H{\o}yer, and Mhalla \cite{DHHM06} show that even deciding if a graph is connected or not requires $\Omega(n^{3/2})$ quantum queries 
in the adjacency matrix model.  
We extend this bound by showing that for any $1 \le \tau \le n$ there is a graph family with edge-weight ratio $\tau$ for which solving the minimum cut problem requires $\Omega(n^{3/2} \sqrt{\tau})$ quantum queries to the adjacency matrix 
(Theorem~35 of the source).  For $\tau \ge n$ one can always use the trivial $O(n^2)$ algorithm, thus 
our results characterize the quantum query complexity of the minimum cut problem in the adjacency matrix model for any value of $\tau$.

For the adjacency array model, we give a bounded-error quantum algorithm that solves the minimum cut problem in an $n$ vertex, $m$ edge graph with edge-weight ratio $\tau$ using $\tO(\sqrt{mn\tau})$ quantum queries (\cref{thm:main}).
The quantum algorithm runs in time $\tO(\sqrt{mn\tau} + n^{3/2})$ (\cref{thm:time}).
In this case we do not know whether the bound is tight in all regimes.
For unweighted graphs ($\tau = 1$) the best lower bound we know of is $\Omega(n)$, which again follows from a lower bound for connectivity \cite{DHHM06}.
For any $\tau > 1$ we show that the minimum cut problem requires $\Omega(n^{3/2})$ quantum queries to the adjacency array (Theorem~37 of the source).
Finally, for any $1 \leq \tau \leq 5n/8$ we show a lower bound of $\Omega(\tau n)$ on the number of quantum adjacency array queries for solving the minimum cut problem (Theorem~40 of the source).

In addition to computing the weight $\lambda(G)$ of a minimum cut, all of our upper and lower bounds also apply to outputting the edges or shores of a cut realizing $\lambda(G)$.


\subsection{Previous work}
We are not aware of any previous work on the quantum complexity of \emph{exact} global minimum cut.  The closest work to ours in topic is the recent paper of Apers and de Wolf \cite{AdW19}, which 
in particular shows that in a weighted graph a $(1+\varepsilon)$-approximation to the weight of a minimum cut can be found in time $\tilde O(n^{3/2}/\epsilon)$ in the adjacency matrix model and time 
$\tilde O(\sqrt{mn}/\epsilon)$ in the adjacency array model.  The sparsifier construction of Apers and de Wolf that yields this approximation also plays a key role in our algorithm.  

Another key work for us is the seminal paper of D\"{u}rr, Heiligman, H{\o}yer and Mhalla \cite{DHHM06} which gives tight bounds for the quantum complexity of many graph problems in both the adjacency matrix and adjacency array models.  In particular, 
they show that determining if a graph is connected or not, i.e.\ determining if the minimum cut value is zero or positive, requires $\Omega(n^{3/2})$ queries in the adjacency matrix 
model and $\Omega(n)$ queries in the adjacency array model.  These are still the best lower bounds we know of for simple graphs\footnote{We use the term simple graph to mean an undirected, unweighted 
graph with no self-loops and no multiple edges.} even for the more general problem of computing the edge connectivity. Indeed, 
we show the $\Omega(n^{3/2})$ connectivity lower bound in the adjacency matrix model is a tight lower bound even on the quantum complexity of edge connectivity.
In \cite{DHHM06} it is also shown that finding a spanning forest in the adjacency matrix model can be done with a quantum algorithm in queries and time $\tO(n^{3/2})$, which is a result we will make use of in our time efficient algorithm.

Two classical papers which inspired our algorithm are the works of Kawarabayashi and Thorup (KT) \cite{KT19} and Rubinstein, Schramm, and Weinberg (RSW) \cite{RSW18}.  
KT give the first near-linear time deterministic algorithm to compute the edge connectivity of a simple graph $G = (V,E)$.  A key idea of KT is to look at a \emph{contraction} of the original graph $G$.  Let 
$\Pcal = \{P_1, \ldots, P_k\}$ be a partition of $V$.  The contraction $G' = \Contract(G,\Pcal)$ is a multi-graph whose vertices are labeled by the sets in $\Pcal$ 
and which has all the edges of $G$ whose endpoints lie in different sets of $\Pcal$.  KT first check the cardinality of all \emph{star} cuts of the form $\Delta_G(\{v\})$, which can be 
done deterministically in linear time.  
To find the minimum non-star cut, KT show that any simple graph $G$ with minimum degree $d$ has a contraction $G'= \Contract(G,\Pcal)$ that 
preserves all of the near-minimum non-star cuts of $G$, but which has only $\tO(n/d)$ vertices and $\tO(n)$ edges.  Moreover, they 
show how to find such a contraction deterministically in near-linear time.  They then use Gabow's $\tilde O(\lambda(G)|E(G)|)$ mincut algorithm \cite{Gabow95} to find a minimum cut in $G'$.
If $G$ has $m$ edges then $\lambda(G') =\lambda(G) \leq m/n$, and as $|E(G')| \in \tilde O(n)$, this gives a time bound that is nearly linear in $m$.

RSW follow a similar high-level approach to give a classical randomized algorithm that computes the edge connectivity of a simple graph with
\emph{cut queries}.  In the cut query model, when the input is a graph $G$, an algorithm can query any nontrivial set $X$ and receive the answer $|\Delta_G(X)|$.  
RSW show that the edge connectivity of a simple graph can be computed with high probability by a randomized algorithm after $O(n \log(n)^3)$ cut queries.  
In fact, this algorithm finds \emph{all} minimum cuts of the graph.  The RSW algorithm again first evaluates all star cuts.
They then
remove the log factors from the KT result to show there is a partition $\Pcal$ of $V$ such that $G' = \Contract(G,\Pcal)$ preserves all 
near-minimum cuts of $G$ and has only $O(n)$ edges.\footnote{An $O(n)$ bound on the number of edges implies an $O(n/d)$ bound on the number of vertices in a black-box way.} 
Moreover, they show how to efficiently learn this contraction with cut queries.  
The log factors of the original KT proof were also removed via another algorithmic proof by Lo, Schmidt, and Thorup \cite{LST20}.  

Our quantum algorithm will follow the approach taken by RSW to learn such a contraction of $G$, as is detailed in the next section.

\subsection{Technical overview}
In this overview we focus on the adjacency matrix model.
Apart from the lower bound, most ideas carry over in a straightforward way to the adjacency array model.  
We start off by explaining the lower bound, as this clearly shows the origin of the~$n^{3/2} \sqrt{\tau}$ complexity.  

\paragraph*{Lower bound on the quantum query complexity.}
For the lower bound we construct a family of graphs on $2n$ vertices with edge weights in $\{1,\tau\}$.  Partition the $2n$ vertices into two sets $A$ and $B$ each of size $n$.  Make 
a complete graph among the vertices in $A$ where every edge has weight $\tau$ and do the same to $B$.  This ensures that $w(\Delta_G(X)) \ge \tau (n-1)$ for any $\emptyset \ne X \subset A$, 
and the same for $B$.  This large value gives us ``cover'' to hide either $k-1$ or $k$ edges of weight 1 between $A$ and $B$.  If $k < \tau(n-1)$ these edges will 
constitute the unique minimum cut, and thus an algorithm that outputs the weight of the minimum cut must determine if we hid $k-1$ or $k$ edges.  This is equivalent to 
determining if there are $k-1$ or $k$ marked items in a search space of size $n^2$, for which a quantum query lower bound of $\Omega(\sqrt{kn^2})$ is known \cite{NW99}.  In our case, with $k = \tau (n-1) -1$ 
this gives a bound of $\Omega(n^{3/2}\sqrt{\tau})$.  Thus we see that ultimately the lower bound for minimum cut boils down to the difficulty of counting for quantum algorithms.  
We will see how a similar task arises in the upper bound as well.  

\paragraph*{Upper bound on the quantum query complexity.}
We first describe a quantum algorithm for computing the edge connectivity of an unweighted graph.  We will follow the outline of the 
RSW cut query algorithm, which proceeds in the following way.  The algorithm first computes the degree of every vertex of $G$, thereby determining the minimum cardinality of a star cut.  The task is then reduced to finding 
the minimum cardinality of a non-star cut.  To do this, the RSW algorithm first produces an $\varepsilon$-\emph{cut sparsifier} of the graph, following an algorithm due to 
Bencz\'{u}r and Karger \cite{BK15}.  An $\varepsilon$-cut sparsifier of $G = (V,E)$ is a sparse weighted graph $H$ whose edge set is a subset of $E$, but where edges 
are allowed to be weighted.  For every nontrivial $X$ the weight of the cut $\Delta_H(X)$ in $H$ is within a factor of $1 \pm \varepsilon$ of $|\Delta_G(X)|$.  

For $\epsilon = 1/100$, the algorithm finds an $\varepsilon$-cut sparsifier $H$ of $G$.  The algorithm is able to write $H$ down in memory and then, without further queries, it can compute the weight of a minimum cut in $H$, say it is $\lambda(H)$, and enumerate all non-star cuts of $H$ whose weight is at most $(1+ 3 \epsilon) \lambda(H)$.  
With high probability this includes the shores of all non-star minimum cuts of $G$.  Let $\Tcal$ be the set of all shores of these cuts.  
The algorithm then computes the coarsest partition $\Pcal = \{P_1, \ldots, P_k\}$ of the vertex set with the property that for all $P_j \in \Pcal$ and $u,v \in P_j$ it holds 
that $u,v \in X$ or $u,v \in \overline{X}$ for all $X \in \Tcal$.  We call $\Pcal$ the set of atoms of $\Tcal$, denoted $\atoms(\Tcal)$. 
As $\Tcal$ is the set of shores of all non-star near-minimum cuts, this means that, for every $P_j \in \Pcal$, no non-star near-minimum cut has an edge with both endpoints in $P_j$;  
as $\Pcal$ is the coarsest partition with this property, among such partitions it minimizes the number of edges \emph{between} components of the partition.  A key 
fact is that $\Contract(G,\Pcal)$ is a sparse graph.

\begin{lemma}[\cite{KT19, RSW18,LST20}]
\label{lem:RSW}
Let $G = (V,E)$ be a simple $n$-vertex graph with minimum degree $d$.  For a nonnegative $\varepsilon < 1$, let 
$\Tcal = \{X : |X|, |\overline{X}| \ge 2 \mbox{ and } |\Delta_G(X)| \le \lambda(G) + \varepsilon d\}$, that is the set of shores of all non-star cuts whose weight is at most $\lambda(G) + \varepsilon d$, and let $G' = \Contract(G, \atoms(\Tcal))$.  Then 
$|E(G')| = O(n)$.
\end{lemma}
By the definition of $\Pcal$ in this lemma, one can also see that $G'$
preserves all of the non-star near-minimum cuts of $G$.  As we already know the minimum degree of $G$, to determine $\lambda(G)$ it suffices to compute 
the edge connectivity of $G'$.  For a query algorithm, to do this it suffices to learn the $O(n)$ edges of the graph $G'$; 
then one can compute the edge connectivity of $G'$ without further queries.  The edge connectivity of $G$ is then the minimum of the minimum degree of $G$ and the 
edge connectivity of $G'$.  

We phrase the RSW algorithm in an abstract way in terms of four computational primitives.  We indicate oracle access to $G$ by square brackets and put the 
parameters explicitly given to the routines in parentheses.
\begin{enumerate}
  \item FindMinStar$[G](\delta)$ --- a routine that given oracle access to $G$ finds the minimum weight of a star cut of $G$ with error probability at most $\delta$.
  \item Cut-Sparsifier$[G](\varepsilon, \delta)$ --- a routine that given oracle access to $G$ outputs an $\varepsilon$-cut sparsifer of $G$ with error probability at most $\delta$.
  \item LearnCutAtoms$(H, \lambda, \delta)$ --- a routine that given an explicit description of a graph $H$, a cut threshold $\lambda$, and an error probability $\delta$, 
  outputs $\Pcal$, the atoms of the shores of all cuts of weight at most $\lambda$, with error probability at most $\delta$.
  \item LearnContraction$[G](\Pcal, M, \delta)$ --- a routine that given oracle access to $G$ and a partition $\Pcal$ of the vertex set, learns $\Contract(G,\Pcal)$ if it has at most $M$ 
  edges and otherwise outputs NULL, again with error probability at most $\delta$.
\end{enumerate}
In \cref{thm:rsw_alg}, we show a general upper bound on the query complexity of edge connectivity in terms of the sum of the query complexity of the routines in steps (1), (2), and (4).  Step~(3) 
requires no queries.  It is somewhat surprising that a randomized algorithm designed for cut queries leads to an optimal quantum query algorithm in the adjacency matrix model.  We hope that phrasing 
the algorithm in this abstract way will make it easy to further apply it to other computational models.

In terms of quantum query complexity in the adjacency matrix model, the cost of the 4 steps are as follows.
Item~(1) can be done with $O(n^{3/2})$ queries by composing the $O(\sqrt{n})$ query quantum minimum finding algorithm over the $n$ vertices with the $n$ query classical algorithm to evaluate the degree of 
a vertex.  The quantum complexity of (2) was recently studied by Apers and de Wolf \cite{AdW19}.  They 
show that even an $\varepsilon$-\textit{spectral} sparsifier can be found in \emph{time} $\tilde O(n^{3/2}/\varepsilon)$ in the adjacency matrix 
model.  For our purposes, we take $\varepsilon=1/100$ giving an $\tO(n^{3/2})$ bound here. Item~(3) costs no queries as the routine is given an explicit description of $H$.  Item~(4) is very similar to the 
problem that we saw in the lower bound: we have to learn up to $M$ edges in a search space of size $O(n^2)$ which can be done with $O(n \sqrt{M})$ queries.  By \cref{lem:RSW}
we can take $M = O(n)$ resulting in an $O(n^{3/2})$ quantum query bound for this step.

These bounds when taken together imply a quantum algorithm for edge connectivity making $\tilde O(n^{3/2})$ queries in the adjacency matrix model.

\paragraph*{Extension to weighted graphs.}
The query complexity of steps~1--3 does not change for weighted graphs.  The complexity of step~4, however, depends on the upper bound $M$ on the number of 
edges in the graph $\Contract(G, \Pcal)$, which does depend on the edge weights.
To extend the above algorithm to weighted graphs, we prove the following generalization of \cref{lem:RSW}.
\begin{restatable}{lemma}{RSWweighted}
\label{thm:RSWweighted}
Let $G = (V,w)$ be a weighted graph with $|V| =n$ and where every edge has weight at most $\tau$.  Let $d = \min_{u \in V} w(\Delta_G(\{u\}))$.  For a nonnegative $\varepsilon < 1$, let 
$\Tcal = \{X : |X|, |\overline{X}| \ge 2 \mbox{ and } w(\Delta_G(X)) \le \lambda(G) + \varepsilon d\}$ and let $G' = \Contract(G, \atoms(\Tcal))$.  Then 
\[
w(E(G')) \le \frac{68 \tau n}{(1-\varepsilon)^2} \enspace .
\]
\end{restatable}
This lemma is tight as can be seen from the cycle graph with all edge weights $\tau$.  Because the bound necessarily depends on $\tau$, applying this lemma 
back to the cut query or sequential models does not seem to lead to good algorithms.\footnote{The randomized cut query complexity of minimum cut for weighted graphs was 
recently resolved using different techniques by Mukhopadhyay and Nanongkai \cite{MN20}.}
For quantum algorithms, however, it is exactly what is needed.

If the edge-weight ratio is $\tau$, for constant $\varepsilon$ \cref{thm:RSWweighted} implies an $O(\tau n)$ upper bound on the number of edges in the contracted graph $\Contract(G,\Pcal)$.
This means that the LearnContraction step can be performed with $O(n^{3/2}\sqrt{\tau})$ queries.  Together with the $\Omega(n^{3/2}\sqrt{\tau})$ query lower bound mentioned 
above we obtain the following tight characterization of 
the query complexity of minimum cut in the adjacency matrix model in terms of the edge-weight ratio.


\begin{restatable}{theorem}{thmquery} \label{thm:restate_main}
Let $G =(V,w)$ be an $n$-vertex weighted graph with edge-weight ratio $\tau$.  There is a quantum algorithm that finds the 
weight and shores of a minimum cut of $G$ with probability at least $3/4$ after 
$\tilde O(n^{3/2}\sqrt{\tau})$ queries to the adjacency matrix of $G$.  Moreover, there is a family of graphs with edge-weight ratio $\tau$ 
for which computing the weight of a minimum cut with bounded-error requires $\Omega(n^{3/2} \sqrt{\tau})$ quantum queries to the adjacency matrix.
\end{restatable}
The upper bound for this theorem is given in \cref{thm:main}, and the lower bound in Theorem~35 of the source.


\paragraph*{Upper bound on the quantum time complexity.}
Let us now consider the time complexity of the above algorithm, corresponding to the total number of queries and elementary gates in the quantum circuit model that the algorithm uses.  Steps~(1) and~(4) are ultimately applications of Grover's algorithm and 
can be implemented in time which is just a $O(\log(n))$ factor more than their query complexity.  For step~(2), Apers and de Wolf already give 
a time complexity upper bound of $\tO(n^{3/2}/\varepsilon)$.  Thus to get an upper bound on the time complexity it suffices to 
analyze the routine LearnCutAtoms$(H,\lambda,\delta)$ from step~(3).
Given a graph $H$, this subroutine requires us to output the atoms of $\Tcal$, where $\Tcal$ is the set of shores of all near-minimum cuts of $H$.  
For this discussion, one should take near-minimum cuts to mean cuts of weight at most $(1+1/100) \lambda(H)$.  
It is known that an $n$-vertex graph $H$ has at most $O(n^2)$ cuts of weight $< 3\lambda(H)/2$ \cite{HW96}.  Thus we know that $|\Tcal|$ is not too large.  
However, we still need to efficiently find these near-minimum cuts.  

To do this we build on Karger's seminal work \cite{Karger00} that connects near-minimum cuts with tree packings.
Consider a spanning tree $T$ of $H$, as in \cref{fig:2-respecting}.
A cut in $H$ with shore $X$ is said to \emph{2-respect} $T$ if it cuts at most 2 edges of $T$, that is $|\Delta_T(X)| \leq 2$.
Karger showed how to efficiently construct a set of $O(\log n)$ spanning trees in $H$ so that every near-minimum cut 2-respects at least one of them.  As each 
tree has at most $n-1 + \binom{n-1}{2} = \binom{n}{2}$ 2-respecting cuts, this family of trees defines a set of shores $\Tcal'$ of cardinality $O(n^2 \log n)$ which necessarily contains $\Tcal$.
A graph can potentially contain $\binom{n}{2}$ minimum cuts, as witnessed by the cycle graph, thus this bound is nearly tight.  
Unfortunately, iterating over $\Tcal'$ is still too costly for us.

\begin{figure}[htb]
\centering
\includegraphics[width=.5\textwidth]{2-respecting.pdf}
\caption{Graph $H$ (thin grey edges) with spanning tree $T$ (thick black edges). The cut with shore $X$ 2-respects $T$ since $|\Delta_T(X)| = |\{e,e'\}| \le 2$.
There are at most $\binom{n}{2}$ such cuts.}
\label{fig:2-respecting}
\end{figure}

As we are only interested in $\atoms(\Tcal)$, and not $\Tcal$ itself, it suffices for us to find a set $\Scal$ such that $\atoms(\Scal) = \atoms(\Tcal)$.  We call such an $\Scal$ a \emph{generating set} for $\atoms(\Tcal)$.  Our next observation is that there necessarily exists a generating set for $\atoms(\Tcal)$ of size $O(n)$.  
This follows by a greedy argument: set $\Scal = \emptyset$ and iterate over all cut shores $X \in \Tcal$, adding $X$ to $\Scal$ iff $\atoms(\Scal \cup X) \ne \atoms(\Scal)$.
The resulting $\Scal$ has the same atoms as $\Tcal$.  Moreover, $|\Scal| \leq n-1$ since every element added to $\Scal$ creates at least one new atom, there are at most $n$ atoms in total, and $\Scal = \emptyset$ has 1 atom.
While a good start, this still leaves the problem of efficiently finding a small generating set.  

We are able to give an explicit description of an $O(n \log(n))$ size generating set.
First consider a single spanning tree $T$ of $H$.
For any $f \in E(T) \cup E(T)^{(2)}$ we let $\shore(f)$ denote the cut shore such that $\Delta_T(\shore(f)) = f$. \footnote{As $\Delta_T(X) = \Delta_T(\overline{X})$, for uniqueness we define a root $r$ in $T$ and choose $\shore(f)$ so that it does not contain $r$.}
Now define an unweighted graph $L(T)$ whose vertex set is $E(T)$ and where $f \in E(T)^{(2)}$ is an edge of $L(T)$ iff $\shore(f)$ is a near-minimum cut of $H$ (i.e., $\shore(f) \in \Tcal$).
We show an example in \cref{fig:line-graph}.
Further, let $O(T) = \{e \in E(T): \exists X \in \Tcal: \Delta_T(X) = \{e\}\}$ index the set of near-minimum cuts that 1-respect $T$.
We prove the following lemma.
\begin{lemma}
Let $\Tcal' = \{X \in \Tcal : |\Delta_T(X)| \le 2\}$ be the shores in $\Tcal$ whose corresponding cuts $2$-respect $T$.  
If $F$ is a spanning forest of $L(T)$ then $\Scal(T) = \{\shore(f) \mid f \in E(F) \cup O(T)\}$ is a generating set for $\atoms(\Tcal')$.
\end{lemma}
Moreover, since $|E(F)| \le n-2$ and $|O(T)| \le n-1$ we have $|\Scal(T)| \le 2n-3$.  Taking the union of $\Scal(T)$ over all of the $\log(n)$ spanning trees $T$ of Karger's tree packing gives a generating set $\Scal$ for $\Tcal$ of size $O(n\log(n))$. 

\begin{figure}[htb]
\centering
\includegraphics[width=.8\textwidth]{line-graph.pdf}
\caption{Left: A spanning tree $T$ (thick black edges) of the graph $H$ (thin grey edges) with minimum cut $\lambda(H) = 2$. Right: The associated graph $L(T)$ with vertex set $E(T)$ and $f \in E(T)^{(2)}$ an edge of $L(T)$ iff $\shore(f)$ is the shore of a near-minimum cut in $H$ (in this case, a near-minimum cut is a cut of weight $\leq \frac{3}{2} \lambda(H)$).}
\label{fig:line-graph}
\end{figure}

We cannot explicitly write down the graph $L(T)$, but using an efficient data structure for evaluating 2-respecting cuts \cite{MN20,GMW20c} we can in $O(\log(n))$ time determine 
whether or not $\{e,e'\}$ is an edge of $L(T)$.
This essentially gives us adjacency matrix access to $L(T)$, and hence we can use the $\tO(n^{3/2})$ time quantum algorithm from~\cite{DHHM06} to construct a spanning forest $F$ of $L(T)$.
We note that it is conceivable that there exists an efficient classical algorithm to do this.
However this would require using further properties of $L(T)$ since classically computing a spanning forest in the adjacency matrix model requires $\Omega(n^2)$ queries.

Once we have the $O(n \log n)$ size generating set $\Scal$, we still cannot naively compute the atoms of $\Scal$ because this would again be too costly.
Rather, we find the atoms of $\Scal$ in $\tO(n)$ time by combining a random hashing scheme with an efficient data structure based on Euler tour trees \cite{HK95}.
This shows that a quantum algorithm can implement step~(3), LearnCutAtoms, in time $\tO(n^{3/2})$.  Note that this running time is independent of the kind of oracle access we have to $G$.  
This gives the following theorem.

